% vim: textwidth=78 ai spell spelllang=en filetype=tex sw=2

\chapter{Introduction}\label{sec:intro}
In the contemporary world, we are surrounded by devices making our lives
easier. It counts from modern lifts which speak to us by natural language to
supercomputers helping to treat cancer such as \emph{IBM
Watson}~\cite{watson-cancer}. Since Moore's law~\cite{moore_law} has proven to
hold for the past four decades, most present devices are featuring quite fast
processing units thanks to the increasing count of built-in transistors. This
allows developers to implement many interesting but complex algorithms.
However being written by human beings vastly, it may lead to errant behaviour
of the implementations as they are only as strong as the weakest link in the
imaginary chain. One example is the error that led to a crash of the
\emph{Ariane~5} rocket that should launch a satellite to the orbit of the
Earth~\cite{ariane_5}. \$370M were buried in the middle of nowhere due to a
programmer's omission to check for an overflow in the acceleration
computation.  So these errors obviously cause companies revenue losses. It is
known the later in the development cycle an error is revealed, the more
expensive it is to fix the error~\cite{topten}. The situation becomes even
worse when the software is deployed at the customer's site or in a rocket
flying 13~000~ft above the sea level already.

Companies have been always desiring for cheap solutions perfectly fitting
their business plans. And yet, solutions which are based on algorithms that
will cause no harm when deployed, completely error-free at best.  But they are
out of luck as, in general, there cannot be a tool which would tell us whether
some algorithm is error-free. If there were one, it would be in contradiction
of the \emph{Halting problem}~\cite{halting_problem}.  However while it is
impossible in theory, it does not necessarily mean we should not at least try
hard for non-general cases when dealing with algorithms implemented in
software. There, the implementations have at least finite state space thanks
to finite memory.  Even though the limited state space is still huge,
companies have been always using techniques helping to avoid and remove errors
from their programs.  Even some kind of \emph{standardization} of used coding
principles like dictated coding style may reduce the number of introduced
errors~\cite{topten}. The reason may be that a project management imposes a
duty on developers to write simpler code. For example, the management can ban
complex constructs which are known to introduce errors. Or, for some
environments like embedded systems, heap allocations can be disallowed. This
results in weeding out one whole category of programming errors -- memory
leaks. Some companies also started adopting standards like
ISO~9001:2000~\cite{iso9000} or MISRA~C~\cite{misra_c}. Their customers are
assured that the internal management of such a company conforms to a
standardized behaviour that tries to increase quality of developed products.

Standardization cannot by definition evade programming errors completely.
Since the early times of computer technology, programmers in companies, or
testers if they were distinct persons, were forced to test the programs.
Testers were running the programs with diverse input, random inclusive, to
make the program misbehave or crash in order to catch and fix program errors.
This approach is called \emph{dynamic program testing}~\cite{sw_testing}. It
is used very often as it is relatively cheap and mostly very effective
technique. Thanks to executing the code in real-time with real data, whenever
the dynamic testing reveals an error, it is guaranteed to be genuine (if we
look aside from malfunction of broken hardware).  But it indeed takes a huge
amount of time to perform systematic testing of large projects. And yet,
dynamic testing needs obviously a runnable executable which is not usually
available in the early phases of development.  Testers can help in that case
by creating simple binaries that exercise only finished subsystems, by
so-called \emph{unit tests}. However tests, especially the manually written
ones, often do not cover all possible paths in the program.  Actually, it was
shown that a certain set of unit tests developed for over fifteen years can be
augmented to cover sixteen more per cents of code. This was achieved by
\klee~\cite{klee} using automatic sophisticated techniques for tests
generation.

\emph{Compilers} can also help to warn about many primitive errors introduced
into a program by a programmer. But compilers are designed to be really fast
in compiling code, so their primary goal is definitely not error finding.
Hence many warnings reported by compilers are rather by-products of performed
optimizations and analyses needed for the intermediate compilation phases. We
have to emphasize that an advanced checking of code got introduced only by
highly optimizing compilers, they were not in the prior compilers at all.
Separate projects like \textsc{Lint}~\cite{lint} or
\textsc{Dave}~\cite{dave} performed this role and were widespread for checking
at that time. These tools were based on pattern matching and also on some
analyses looking at how data flow in the checked code. Mostly, these
techniques were more computationally expensive than compilers, but for that
price, the tools implementing them were giving more accurate results with more
reports overall.  These tools were the first attempts to use the \emph{static
program analysis}~\cite{program_analysis}. The static analysis tries to solve
the problems of finding errors in programs automatically, possibly in the
early phases of development, and more importantly in a whole program and
without the urge of the dynamic execution. But unlike the dynamic execution,
static techniques are often prone to the presence of reports that are not
real. Those are errors found in the checked code by an analyzer, but they can
never occur in practice, only the analyzer is imprecise. There are many ways
to cope with such reports.  For example, tools like \textsc{Lint} suggest a
programmer to annotate the checked code, so that the tools have a clue and can
thus generate less invalid reports.

%\todo{overview in~\cite{analysis_old_survey}}
In fact, much more sophisticated techniques like of Cousots~\cite{cousot} or
King~\cite{se_king} were described theoretically couple years before
\textsc{Lint} was presented.  However they came of real use even lately, when
computers become powerful enough to finish computation of such analyses in
reasonable time.  Now, having the computation power and multi-processor
(thousand cores are not exceptions today) machines, we witness a rapid
improvement of the situation over the past two decades. We could also see many
new techniques emerging. Many of those were adopted or even developed by
industrial companies to make their products safer. For example, Microsoft
introduced SDV toolkit~\cite{sdv} based on \textsc{Slam}~\cite{slam_cav,
slam_popl, slam_decade} and later SDV 2.0 with \textsc{Slam2}~\cite{slam2} to
improve situation around Windows drivers.  Microsoft also released for
internal use \textsc{Sage}~\cite{sage} which found errors in core Windows
parts like the parser of cursor files. Other companies like Coverity, Gimpel
Software, GrammaTech, Klocwork, or MathWorks presented their
products~\cite{coverity, pclint, codesonar, klocwork, polyspace} to offer a
service to check code of their customers. Some of the companies also check
open source code like the Linux kernel or core Linux libraries. Although the
companies make the results available~\cite{coverity_scan}, developers of the
kernel or libraries usually have no control over the setup and check
frequency, therefore they have no chance to check their code on demand.

For those reasons, some groups decided to implement their own, freely
available, checkers. For example \textsc{Cppcheck}~\cite{cppcheck},
\fbugs~\cite{findbugs}, \smatch~\cite{smatch}, \sparse~\cite{sparse}, or
\stan~\cite{stanse}.  They vary in precision, speed, and of course in the used
method for finding bugs.  Although developers may use them at any time at
their will, they often get results with much higher unreal reports rate than
results of a commercial tool.

\section{Motivation}
In an ideal world, every code checking tool is capable to handle and check any
project. Unfortunately, this is not true in the real world we live in. What
are the difficulties we can hit when analyzing large software projects? And
why do we should care? Those are exactly the questions which this section
gives answers for. First, specifics of large projects are enumerated. Later,
we explain why it is important to find as many errors as possible in them.

\subsection{Requirements for Analyzers}
If code analysis tools were perfect, we could probably skip this section.
But this is not true, the tools are often incomplete in some sense when it
comes to analyzing real large-size code.  We can likely conclude that every
problem in handling of large software projects stems from a single source:
\emph{complex code}. Now let us split the source into several categories and
comment each of those. An observing reader may notice that not all of the
categories are specific solely to large projects. But we understand the list
rather as an enumeration of items which are a must for a code analyzer
evaluating large projects.
%This will happen from the most abstract level to simple code constructs.

\begin{description}
\item[Organization of Files] Source code with project's functionality is often
divided into files and files are organized in directories. It is important to
know where the analyzer can find all those files which constitute the project.
Should a checker just recursively crawl a specified directory and check all
sources in there?  What if some source file includes another one and the inner
is not compilable on its own as happens in the Linux kernel?  Even if a user
would specify the files manually, large projects tend to pass special
directives to a compiler. The directives often pass paths to header files,
definitions of some macros from the command line, or pre-inclusion of specific
files.  And since the user of the analysis tool need not be a developer of the
project, writing the options down by hand is infeasible for large projects.
Checking tools which work on the source level definitely need to support
\emph{build systems} and infer the files and compiler options automatically.
They cannot depend on the user's input.

\item[Source Code Language] It is not often easy to read and parse source
code of a large project. This is due to languages used for writing code of
such projects.  Even though there are standards for languages, there are many
language offshoots that ``evolved'' with respective compilers. Let us consider
C with the last standard issued in 1999~\cite{ansi_c}. There are many
deviations from the standard introduced by \gcc~\cite{gcc}, \textsc{Intel
CC}~\cite{intel_cc}, \textsc{MSVC}~\cite{msvc}, and possibly others.  Every
single of them supports a different superset of the standard.  Thus it is very
hard for analyzers to implement a parser which would be universally capable of
parsing whichever program with sources in a specified language. This causes
problems when analyzing large projects since there is a high probability of
using many features of some offshoot implemented in some default used
compiler.
%code written by diverse people+large chance to hit something unusual in the
%large code: all diverse kind of code, whatever one can think of

\item[Much Parsed Information] When a checking tool is able to parse some
project's code, intermediate code structures are created from the code. But
information about all these structures barely fits into the memory for large
projects. We have to think how and which information to store and where. For
example dumping the information into files and restoring them later. Or
computing only minimal data which sum the effect of the parsed code and
dropping the large structures afterwards.

\item[High Computation Cost] Some techniques cannot be applied to real
software due to a high computation cost.  When developers of large project
analyzers choose techniques to implement in their analyzer, they must have in
mind the large amount of code which needs to be handled. The state space to
be explored can be larger by many orders of magnitude in comparison to an
analysis of small programs. Although computers are fast enough to deal with
more precise analyses today, we are still not there to fully prove there is no
error in some large project. The increase of computation power can be
demonstrated on the symbolic execution~\cite{se_king}. It started being used
almost thirty years after its introduction because the computers were not so
powerful and nobody could really use it for anything real.

\item[Higher Level Analysis] The best code practices suggest to split
functionality among smaller logical pieces for better readability and
maintainability. Programmers usually put these pieces into separate functions.
But these pieces alone are not useful when finding errors. They are potential
sources of invalid reports. Consider the example of
Figure~\ref{fig:intro_need_context}.
\begin{figure}[htb]
\begin{center}
\begin{tabular}{c}
\begin{lstlisting}
DEFINE_LOCK(some_lock);

void my_lock() {
  lock(&some_lock);
}
void my_unlock() {
  unlock(&some_lock);
}
void fun() {
  my_lock();
  my_unlock();
}
\end{lstlisting}
\end{tabular}
\end{center}
\caption{Code to demonstrate the need of analyses to work on a higher level
than functions.}
\label{fig:intro_need_context}
\end{figure}
There, \code{my_lock} locks the lock \code{some_lock}, \code{my_unlock}
unlocks it and they are called in a sequence from \code{fun}.
If we analyzed \code{my_lock} and \code{my_unlock} alone, we would receive two
invalid reports. First, that \code{my_lock} is left with \code{some_lock}
still held. Second, unlock of an unlocked lock in \code{my_unlock}. But when
all three are analyzed together, no such report is generated. This drags in
the need of more complex analyses which look at code from a higher level.

\item[External Dependencies] Large projects tend to use some external library
very heavily, e.g.\ for reusability purposes. So to have a complete analysis
these libraries have to be analyzed too. But there is a slight problem when
the libraries are written in a different language. The analyzer still can be
taught to understand that. But it can be worse, when no sources are available
at all.  There are still some possible solutions. First, we can develop an
analyzer working on the binary level where no sources are needed. This was
demonstrated in \textsc{Sage}~\cite{sage} for example. Second, we can ignore
calls to unknown code. But we lose precision of course. Finally, we can solely
aim at projects that are self-contained. They do not depend on any other code.
These are kernels of operating systems for example.

\item[Wide Range of Properties] Large projects usually contain diverse parts,
which contain a handful variety of errors. Not necessarily, but analyzers are
desired to report as many various reports as possible during a single run.
Even if a project is highly specialized, i.e.\ there is some class of errors
missing, the tools still should find some other classes of errors. Running
several tools is possible and often performed, but it is much easier to run a
single tool which covers more error types.

\item[No Manual Annotations] Annotations are used sometimes to suppress
invalid reports from the output. They give clues to the analyzer about the
code which is analyzed. But unlike code of smaller projects, code of large
projects cannot be annotated manually. First, because nobody really knows all
interactions fully as the projects are usually developed by independent groups
of people.  Second, it is not trivial to annotate such a large project
manually. The burden posed on developers would be unbearable.

\item[Concurrency] It is of high probability, that large projects support
creation of threads which can run in parallel. This adds another class of
errors to them: concurrency errors. These are errors that result from an
unexpected scheduling of threads, mostly because invalid or missing critical
sections (locking).  Many checked properties like invalid pointer dereferences
do not necessarily need to be thread-aware. But tracking complex locking
errors or deadlocks using code analysis needs an analyser, that can handle
threads creation and scheduling of processes.

\item[Complex Environment] Usually, implementations of projects run in an
environment that is somehow initialized. It is a running kernel which serves
system calls, disk content, input of a user, etc. The larger a project, the
more complex environment. Analyzers have to cope with that, usually by
modelling the environment. But it is not an easy task.  If modelled
incorrectly, especially precise analyzers relying on the precision would
produce invalid reports.

\end{description}

There are perhaps more requirements like those listed above, we tried to list
especially the important ones.  We can sum up with an overall view: the
analyzers have to be as much precise as possible, without much user
intervention, but still reasonably fast to finish the analysis in the
foreseeable future. A research and implementation of such an analyzer is
indeed a challenge!

\subsection{Picking the Project}
Since this thesis wants to deal with analysis of large projects, we have to
select some project for those purposes.  There exist hundreds of software
projects of large size which we could consider. Let us first divide the
projects roughly into two categories.  According to the design of operating
systems there are two basic privilege levels in which programs run~\cite[pages
21--23]{os_concepts}.  Kernels of the operating systems often run at the level
with no limitations, i.e.\ in the \emph{kernel mode}.  There, the kernel can
touch any memory in a machine and can also change other possible reachable
resources.  On the other hand, user space programs tend to be restricted,
i.e.~they run in the \emph{user mode}. There, a process cannot change any
system resource or memory of others, if not explicitly allowed to, e.g.~by
shared memory~\cite[pages 124--126]{os_concepts} or special system
calls~\cite[pages 62--73]{os_concepts}.  In particular, whether some code runs
in the kernel or user mode has a direct impact on the effect caused by a
potential error in the code. A crash of the kernel means that the machine
ceases production any further actions, whereas a crash of a user space
application causes mostly only pain to its user.

At the first glance it looks like we should specifically aim at kernels of
operating systems. But think about a vulnerability in a random web browser or
database system.  It could indeed lead to stolen user's personal data or even
money from their bank accounts. The same as a problem in the kernel. And such
projects consist of millions lines of code too. But operating system kernels
should generally be error-free more than user space programs, because even a
minor error would allow an arbitrary attacker to obtain administrator
privileges on a machine running such a vulnerable system, or simply crash it.
For a \emph{denial of service} (DoS)~\cite{dos} caused by a crash, it
sometimes suffices to write a ``killer'' string to a special file locally or
even send a packet with a specific structure over the network. Note the
attacker can do it remotely from a place thousands miles away.  Instances of
such errors were already reported to the Linux kernel
developers~\cite{exploit_vmsplice, exploit_proto_ops, exploit_userns}, in a
book about hacking~\cite[page 256]{hacking}, and more to be found in the
National Vulnerability Database~\cite{nvd}. If the attacker succeeds beyond
crashing a machine and is able to get a login there, they are very close to
stealing whatever data they want from that machine since it is much easier to
gain an administrator shell when logged on as a casual user already (sometimes
simply because remote administrator logins are not allowed). And yet it is
often easier to gain an access to other machines in the same subnetwork than
attacking them from the outside. One of the reasons is that there are usually
no managed firewalls for the traffic flowing among computers inside the
subnetwork~\cite{insider_attack}.

The arguments stated above persuaded us to aim at operating system kernels.
But we can choose from many. There are kernels of BSDs (NetBSD, FreeBSD,
OpenBSD, \dots), GNU Hurd, GNU Mach, Linux, Mac OS, Minix, NetWare, OS/2, Sun
OS, Solaris, Windows, and many others. Nonetheless, we limited ourselves to
the Linux kernel. The main reasons are that the kernel:
\begin{itemize}
\item is an open source,

\item is a large code base consisting of more than fifteen millions lines of
code (version 3.10),

\item is developed so fast that each version released every two--three months
contains over 700'000 changed lines of code which are across more than eleven
thousand files (averaged for all releases between 3.5 and 3.10), and

\item is used widely, reportedly on 42\,\% of consumer computing
devices in 2012~\cite{linux_42_percent}.

\end{itemize}
Let us emphasize that the Linux kernel, the same as other kernels, has another
positive aspect: it is self-contained.  And there is still the author of this
thesis who is always biased when choosing a project for code analysis. He
constantly spends some time by Linux kernel development, so he knows the
processes and has some considerable knowledge about the code too.

\section{Background}
This section introduces the reader into the terms we use throughout the
thesis. First, we introduce important terms that describe usefulness of
reports. Next, since every implemented analysis has some properties, we
describe the properties which we consider important. Such properties are
usually looked at when deciding which of the available implementations to
choose for a newly developed program, a static analyzer in particular.  There
is also a small introduction to code parsing as we deal with source codes in
this thesis and they need to be parsed into internal structures.

\subsection{Classes of Reports}\label{sec:intro_reports_cls}
While code analyzers are finding errors in some project, they produce reports.
Not all the reports are useful. Some of them tend to be invalid with respect
to the specification. Imprecise analyses can report errors that are found in
analyzed code, but will never occur when the code is really executed. The
reports have to be somehow classified into the classes which we describe
below.

\begin{itemize}
\item \textbf{Errors} -- This is obviously a correct report. When a tool
reports it and we conclude it can really happen in the analyzed code, we
classify it as an error.

\item \textbf{False reports} -- This includes two classes which we describe
below. Both of them are invalid and should be minimized in real code
analyzers as they bother users. When we mean both of them, we use this term.

\begin{itemize}
\item \textbf{False positives} -- This is an error reported by an analyzer,
which can never happen when the analyzed code is really executed.

\item \textbf{False negatives} -- These are only imaginary reports as they do
not exist. These are errors present in code which an analyzer left unnoticed,
i.e.\ it missed real errors.

\end{itemize}

\end{itemize}

\subsection{Theoretical Properties}\label{sec:theor_props}
Usually, the most important aspect for users of an analyzer is to know whether
the produced output of the analyzer makes some sense according to what is
expected, i.e.\ if the results are \emph{correct} (sometimes called
\emph{sound}). Next, we usually care whether all possible errors are reported,
i.e.\ whether the output is \emph{complete}. And perhaps if there will be any
output in finite time at all, if the analyzer will \emph{terminate}. Below, we
took the definitions over from Cousot et al.~\cite{soundness} and modified
them for our needs. Furthermore, we provide examples of soundness and
completeness in Table~\ref{fig:theor_props} for better understanding.

\begin{description}
\item[Soundness]
A program analyzer is sound when it can \emph{never} omit to signal an error
that can appear at runtime in some execution environment. Sound analyzers are
error eradicators since they find all bugs in a well-defined class (that of
runtime errors). Unsound analyzers are only error finders aiming at finding
some of the errors in the well-defined class. Their main defect is
unreliability, being subject to false negatives thus claiming that they can no
longer find any error while many may be left in the considered class.

\item[Completeness]
A static analyzer is complete when it \emph{never} produces false reports on
any program. This terminology is justified by the fact that such an analyzer,
when queried whether or not some true property holds on the program (for
instance: ``the program never ends with an overflow''), will never issue an
alarm. Never producing any report, or only a few definite ones, would be
complete (but useless). So this completeness requirement should be understood
in the context of soundness.

\item[Termination]
Some static analyzers are subject to non termination, running forever until
exhaustion of resources. This is the case of model checking for
example~\cite{soundness}.

\end{description}

\begin{table}[tb]
\begin{center}
\begin{tabular}{|l||c|c|c|} \hline
\multicolumn{1}{|l||}{\multirow{2}{*}{\textbf{Code}}} &
  \multicolumn{3}{c|}{\textbf{Value Analysis of \code{a}}} \\ \cline{2-4}
  & \textbf{Sound} & \textbf{Complete} & \textbf{Sound and Complete} \\ \hline\hline
\code{a = 1;}    & 1           & 1        & 1    \\ \hline
\code{if (x)}    & 1           & 1        & 1    \\ \hline
~~\code{a = 2;}  & 2           & 2        & 2    \\ \hline
\code{return a;} & $\emptyset$ & 1, 2     & $((x = 0) \Rightarrow 1) \vee ((x
					\neq 0) \Rightarrow 2)$  \\ \hline
\end{tabular}
\end{center}
    \caption{Example to demonstrate soundness and completeness of analyzers.}
  \label{fig:theor_props}
\end{table}

In an ideal world, we would have sound and complete analyzers. They would
return only results without false positives and false negatives. But
computation time of such analyzers tend to be high as well as memory
consumption. This stems from the need of too much precise computation and
precise results which would have to be stored somewhere. So it is very common
to implement analyses that are only complete, but unsound. They report
everything which is possibly an error, but sometimes, we have to await
false positives. There also exist analyzers which are only sound, but
incomplete, or the third option: analyzers that are neither sound, nor
complete. The latter may help in reducing both false positives and false
negatives for the price of potentially hiding errors, i.e.~adding
false negatives.

\subsection{Sensitivity of Analyzers}\label{sec:sensitivity}
Another interesting property of many analyzers is how they are handling
separate commands and functions present in some code. Consider the example of
Figure~\ref{fig:lst_sens} and also consider an analysis that is watching
values stored into variables and is thus able to report this to the user.
\begin{figure}[htb]
\begin{center}
\begin{tabular}{c}
\begin{lstlisting}[language=C]
int fun(int a) {
  int x = 1;

  if (a == 2) (*@ \label{lst:sens_cond} @*)
    x = a; (*@ \label{lst:sens_xa} @*)
  x = 3; (*@ \label{lst:sens_x3} @*)
  return x; (*@ \label{lst:sens_ret} @*)
}
void fun2(void) {
  fun(5);
}
\end{lstlisting}
\end{tabular}
\end{center}
\caption{Code used for demonstration of analyses sensitivity.}
\label{fig:lst_sens}
\end{figure}
Given the example, we describe the following three basic properties which we
usually consider when choosing an analyzer:
\begin{itemize}
\item When an analyzer analyzes some code, it may take into account the order
of program statements and remember what is followed by what.  For the code in
our example this would mean the analyzer will report that the only possible
value in \code{x} at line~\ref{lst:sens_ret} is simply 3 (the last store to
\code{x} before \code{return} proper). The other option is to completely
ignore the order and with considerably less precision report that \code{x} at
line~\ref{lst:sens_ret} can be any of 1, value of \code{a}, or 3 (the analyzer
is not aware of what precedes the \code{return}, hence it has to report all
the sets).  This is called \emph{flow sensitivity} or \emph{insensitivity},
respectively. We often speak about this property when dealing with one
particular function.

\item An analyzer may or may not remember conditions used for branching. If it
does, in our example, \code{x} is then reported to be one of 1, 2, or 3,
because \code{a} was noted to be 2 before the assignment on
line~\ref{lst:sens_xa}. In the other case, \code{x} is announced to be 1, 3,
or whatever value the analyzer previously computed \code{a} can be. We have
just described \emph{path sensitivity}. Obviously, it subsumes flow
sensitivity.

\item An analyzer is called to be \emph{context-sensitive}, or
\emph{inter-procedural} if it considers a calling context (call-stack).  In
our example, a context-sensitive analyzer would say \code{x} can be 1 or 3
since the condition \code{a == 2} is false on line~\ref{lst:sens_cond}. The
counterpart would say \code{x} can be any value due to the assignment on
line~\ref{lst:sens_xa} (the analyzer has to assume \code{a} can contain
whatever value to remain sound).
\end{itemize}

\subsection{Types of Analyses}\label{sec:types_anal}
Analyses are also split into sets of what information they provide. There
are two properties in this category presented below.

\begin{description}
\item[May/must] Sometimes we are interested whether some result given by an
analysis must hold definitely or may hold sometimes. If we get back to the
example of Figure~\ref{fig:lst_sens}, when an analysis says \code{x}
\emph{must} be a constant at line~\ref{lst:sens_ret}, we can just return that
constant and eliminate all the preceding code. This would not be possible with
an answer ``maybe'' given by the \emph{may} analysis. Nonetheless, may
analysis is useful for example if we are interested whether a divisor in a
division may be zero sometimes, in which case we would report an error.  Must
analysis does not necessarily find this bug.

\item[Backward/forward] For flow-sensitive analyses, the way of propagation of
the information is important. The complexity of the analyzer usually goes hand
in hand with the direction used. For example, we may want to know at each line
of a program what variables are used further in the code (to see whether we
need to hold their values still in registers for example). This information is
often propagated backwards from the end to the beginning of a function where
the variable is defined.  Contrary, to see at each line which variables are
already set to some value, analyzers usually propagate the information
forwards.
\end{description}

\todo{emphasize that this thesis deals with practical aspects of using static
analyzers}

\subsection{Other Views}
Of course the list above is not complete. There are many other categorizations
we do not list in here. In short, we can mention white-box versus black-box
where the algorithm either knows the internals of the checked projects or not.
For black-box, an analyzer just can start finding bugs, but will most likely
cover less code than with a white-box analyzer that has to find some knowledge
about the code first.  An analogy might be seen in checking a engine prior
buying a car. Asking a seller for taking a ride, check whatever we see and
check whether an output from that is as expected (speedup after pressing the
throttle pedal for illustration) is black-box because we are unaware about the
internals. Asking a technician to inspect the engine of the car is white-box,
because he has the knowledge about the internals. There is a mid-way called
grey-box testing. The technique requires only an overview of the checked code,
but not full understanding of what is going on in there.

Another angle of view might be hunting functional versus performance issues.
Once a project is working nicely with respect to the demanded functionality,
we may still want to look where we can speed the implementation up. This is
often seen in operating system kernels development as behaviour of a kernel
has to be transparent and as fast as possible, the effects of the kernel
should become invisible at best.

\subsection{Code Parsing}
We presented the terms above directly on text written in programming
languages. This is a form understandable by people when reading, but it is not
so easy to work with when writing an analyzer. Analyzers are used to work on
more convenient structures representing the source code.  There are several
structures used in code analysis and compilation.  This section will cover the
process of creating these structures from the human readable format, including
a description of the created structures proper.

First, to be able to understand source code on input, we must know how to
interpret it. For each programming language, we have a set of rules. They
define allowed \emph{characters} (usually basic alphabet with parenthesis,
curly braces, and other symbols), how to build words from them and their
meaning: what are identifiers, reserved words etc. Finally, the rules specify
how a well-defined sentence looks like, i.e.\ how to build expressions and
statements, functions and whole files. This is accompanied with an explanation
what each such block means, that is \emph{semantics}. So when we write for
example ``\code{char X;}'' in C, we know we have declared a variable which is
capable to hold at least 8 bits (an ASCII character).  The rules are defined
for most languages in the form of a \emph{standard}. For instance, for both C
and C++, ISO standards exist~\cite{ansi_c, ansi_cpp}, whereas for Java, the
rules are defined by the Sun's implementation. Despite Sun tried to
standardize Java twice, they withdrew the drafts both
times~\cite{java_non_std}. So Java is a \emph{de facto} standard defined by
the Sun's implementation.  Remember that even if a standard exists, there is
no force to push compiler implementors to adhere to that standard. This is
seen in the C family and its compilers. What is compilable by GNU's C
compiler~\cite{gcc} is not necessarily understood by the Intel's
compiler~\cite{intel_cc} for example.

\begin{figure}[tb]
\begin{center}
    \tikzstyle{pre} = [<-,shorten <=1pt,>=stealth',semithick]
    \begin{tikzpicture}[node distance=5mm]
        \node [tl_state] (S) {Source};
        \node [tl_phase] (LEX) [right=of S] {Lexer}
          edge [pre] node {} (S);
        \node [tl_state] (TOK) [right=of LEX] {Tokens}
          edge [pre] node {} (LEX);
        \node [tl_phase] (PAR) [right=of TOK] {Parser}
          edge [pre] node {} (TOK);
        \node [tl_state] (AST) [right=of PAR] {AST}
          edge [pre] node {} (PAR);
        \node [tl_state] (CFG) [above right=of AST] {CFG}
          edge [pre,dashed] node {} (AST);
        \node [tl_state] (OTH) [below right=of AST] {\dots}
          edge [pre,dashed] node {} (AST);
    \end{tikzpicture}
\end{center}
    \caption{Process of creating program structures which are more convenient
    to work with than source codes.}
  \label{fig:parsing_process}
\end{figure}

With the rules, we can start writing a translator that reads source code in
the human readable form on the input and produces some structures on the
output. Such translators are called \emph{parsers}.  The whole process of
parsing is depicted in Figure~\ref{fig:parsing_process}.  As we can see, there
are several stages. We discuss them in further text one by one.  The
explanation happens on the C programming language, more concretetely on the
code of Figure~\ref{fig:intro_parsing_codes}~(a).

\begin{figure}[htb]
\begin{center}
\begin{tabular}{c}
    \begin{tabular}{ccc}
    \begin{lstlisting}[language=C]
void fun(int a)
{ /* comment */
  int b = 0;
  if (a)
    b = 5;
  a = 3 * b;
}
    \end{lstlisting} & & {
    \begin{lstlisting}[numbers=none]
key:void id:fun LPAR key:int id:a RPAR
LCUR
  key:int id:b EQ num:0 SEMIC
  key:if LPAR id:a RPAR
    id:b EQ num:5 SEMIC
  id:a EQ num:3 MULT id:b SEMIC
RCUR
    \end{lstlisting} } \\
    (a) & & (b)
    \end{tabular}
\end{tabular}
\end{center}
    \caption{(a) code example for parsing demonstration; (b) tokenized code.}
  \label{fig:intro_parsing_codes}
\end{figure}

%Some standards also contain a \emph{formal grammar}.

First, the input is read by \emph{lexical analyzer} (lexer) character by
character. It consults with language rules whether the input is well-formed
and it starts creating so-called \emph{tokens}. Tokens are character sequences
that are marked by information whether they are a string, number, reserved
keyword (\code{int}, \code{const}, \dots), identifier (\code{some_variable}),
or any other kind defined in the rules. We can see how our example was
tokenized in Figure~\ref{fig:intro_parsing_codes}~(b).  Keywords, identifiers
and numbers are prefixed by type, braces and parenthesis are converted to
special names, the same as for operators.

A list of these tokens is then passed to the input of an \emph{syntactic
analyzer}. Again, it implements the language rules and looks whether the token
sequences are correct. In case they are, syntactic analyzers usually build an
\emph{abstract syntax tree} (AST). All irrelevant information is forgotten
there (abstracted from the original language) and the rest is stored in a tree
as is depicted in Figure~\ref{fig:intro_parsing_ast}.
\begin{figure}[htb]
\begin{center}
\begin{tabular}{c}
    \begin{tikzpicture}[level/.style={sibling distance = 2.5cm/#1,
    level distance = 1cm}]
      \node [ast] {\texttt{fun}}
	  child [sibling distance=1.8cm] { node [ast] {\texttt{void}} }
	  child { node [ast_hl] {param}
	    child { node [ast] {\texttt{int}} }
	    child { node [ast] {\texttt{a}} }
	  }
	  child [sibling distance=4.5cm] { node [ast_hl] {body}
	    child [sibling distance=2cm] { node [ast_hl] {decl}
	      child [sibling distance=1.3cm] { node [ast] {\texttt{int}} }
	      child { node [ast_hl] (L_1) {var}
		child { node [ast] {\texttt{b}} }
		child { node [ast] {\texttt{0}} }
	      }
	    }
	    child [sibling distance=2cm] { node [ast] {\texttt{if}}
	      child { node [ast] (L_2) {\texttt{a}} }
	      child { node [ast] (L_3) {\texttt{=}}
		child { node [ast] {\texttt{b}} }
		child { node [ast] {\texttt{5}} }
	      }
	    }
	    child [sibling distance=2cm] { node [ast] (L_4) {\texttt{=}}
	      child { node [ast] {\texttt{a}} }
	      child { node [ast] {\texttt{*}}
		child { node [ast] {\texttt{3}} }
		child { node [ast] {\texttt{b}} }
	      }
	    }
	  }
	  ;
	  \node [red,font=\scriptsize,above=-1mm of L_1] {1};
	  \node [red,font=\scriptsize,above=-1mm of L_2] {2};
	  \node [red,font=\scriptsize,above=-1mm of L_3] {3};
	  \node [red,font=\scriptsize,above=-1mm of L_4] {4};
    \end{tikzpicture}
\end{tabular}
\end{center}
    \caption{Abstract syntax tree of the code of
    Figure~\ref{fig:intro_parsing_codes}.}
  \label{fig:intro_parsing_ast}
\end{figure}
At this abstract point we have no clue what the original language was.
Compilers usually have separate modules for different languages, called
\emph{frontends}. They parse the input and create one pre-defined form of the
AST. Every frontened does so, hence we have one type of the AST for all
diverse languages the frontends support.  Then, the rest of the compiler
pipeline can be the same for all of the languages and can run the same
optimization passes for them all.  The tree is held in the compiler through
the whole compilation process. It is usual that further passes in the pipeline
add more information to the nodes of this tree. The very common one is a
position of the nodes in the source code for debugging purposes. Even though
ASTs usually hold information about the whole parsed input, we omitted for
simplicity other functions, global variable declarations and other constructs
from Figure~\ref{fig:intro_parsing_ast}.  There is only a subtree
corresponding to the function from Figure~\ref{fig:intro_parsing_codes}~(a).

From the AST, we can derive much information about the parsed code. One use is
to build a \emph{control flow graph} (CFG). CFG of a function is a directed
graph where nodes contain statements of that function and edges are labeled
by conditions when coming from branching statements. In a CFG, the \emph{entry
point} of a function is a successor of a triangle standing on a vertex.  The
\emph{exit point} (return from a function) is a predecessor of a triangle
standing on an edge.  A \emph{branch} is a node with more than one successor.
As an example we have built a CFG of the \code{fun} function in
Figure~\ref{fig:intro_parsing_cfg}. Nodes of the CFG refer back to the to AST
with numbers to see where the statements come from (every node or subtree in
an AST have its rule to create nodes in a CFG).
\begin{figure}[htb]
\begin{center}
\begin{tabular}{c}
    \begin{tikzpicture}[node distance=5mm]
      \node [cfg_start] (S) {};
      \node [cfg] (L_1) [label=right:\red{1},below=3mm of S] {\code{b=0}}
        edge[cfg_edge] node {} (S);
      \node [cfg] (L_2) [label=right:\red{2},below=of L_1] {\code{a}}
        edge[cfg_edge] node {} (L_1);
      \node [cfg] (L_3) [label=left:\red{3},below left=of L_2] {\code{b=5}}
        edge[cfg_edge] node [label=above left:\code{a}] {} (L_2);
      \node [cfg] (L_4) [label=right:\red{4},below=1.3cm of L_2]
	{\code{a=3*b}}
        edge[cfg_edge] node [label=right:\code{!a}] {} (L_2)
        edge[cfg_edge] node {} (L_3);
      \node [cfg_end,below=3mm of L_4] {}
        edge[cfg_edge] node {} (L_4);
    \end{tikzpicture}
\end{tabular}
\end{center}
    \caption{Control flow graph of the code of
    Figure~\ref{fig:intro_parsing_codes}.}
  \label{fig:intro_parsing_cfg}
\end{figure}
Alternatively, for easier explanation of some algorithms, statements can be
put on edges instead of putting them into nodes.  The control flow graph is
not the only structure that can be derived from the abstract syntax tree. We
can derive symbol tables, type definitions, perform semantic analysis to
verify more than syntactic rules etc.

Even though, it is not easy to write a well-working parser by hand, some
projects like \gcc do so. But there also exist tools that help to generate
lexical and syntactic analyzers automatically from a specified formal grammar.
Depending on the programmer's preference on the type of the grammar, there are
various tools. For LL(*) grammars~\cite[Chapter 4.4]{dragon}, one of the most
known tools is \textsc{Antlr}~\cite{antlr4}. \textsc{Flex} and
\textsc{Bison}~\cite{flex_bison} are used to generate LR parsers~\cite[Chapter
4.5]{dragon}. This introduction to parsing does not want to pretend it is
complete, but we do not need more information for the purposes of this thesis.
More information on parsers can be learnt from various literature about
compilers~\cite{dragon, muchnick}.

\subsection{Pointer Analysis}\label{sec:pointer_analysis}
For languages using pointers like C, it is important to know where these
pointers can point to. Consider the global integer variable and an integer
pointer as are in the code depicted in Figure~\ref{fig:lst_ptsto}.
\begin{figure}[htb]
  \begin{center}
  \begin{tabular}{c}
  \begin{lstlisting}[language=C]
int fun(int a) {
  static int x;
  int *p = &x;
  *p = a;
  return x;
}
  \end{lstlisting}
  \end{tabular}
  \end{center}
  \caption{Pointer analysis demonstrative example.}
  \label{fig:lst_ptsto}
\end{figure}
Without knowing that \code{p} can point to and actually changes \code{x}, we
would conclude that \code{fun} always returns 0. But it returns whatever it
gets on the input as the parameter \code{a}. To cope with this situation,
algorithms known as \emph{pointer analyses} exist.  There are two basic
approaches to pointer analysis widely used nowadays:
Steensgard's~\cite{steensgaard} and Andersen's~\cite{andersen}. Both analyses
are \emph{may}. Nonetheless, while the former is based on equalities of the
sets of possible pointer targets, the other on their subsets. So the former
makes the sets of targets equal whenever two pointers can point to a common
target variable.  The latter only copies some (or all) members. In general,
this makes subset-based analyses more precise, but slower. There is an
approach that tries to speed up the Andersen's analysis while preserving its
accuracy.  This approach was suggested by Shapiro and
Horwitz~\cite{shapiro_horwitz}.  They divide the program variables into
categories. Then, it works similar to Steensgaard's algorithm, except they
make target sets equal only if they are in the same category. The analysis
is then repeated with information about categories from the previous run until
it stabilizes. This often results in more accurate results than Steensgaard's
algorithm, but it is faster than Andersen's. The speed of the Shapiro-Horwitz
approach is still lower than of Steensgaard, nonetheless.

The \emph{alias} or sometimes \emph{points-to analysis}~\cite{alias_anal} is a
bit different to the pointer analysis in what it is able to tell about
pointers.  The alias analysis can only conclude that two pointers either point
to the same location or not. It cannot report the set of possible targets.
This is sometimes sufficient and for example the compiler framework \llvm as
of 3.1 provides its optimization passes only with the alias analysis. Precise
alias analysis is an undecidable problem, often a \emph{may-alias}
over-approximation or \emph{must-alias} under-approximation is computed
instead. Another view to pointers suggests the \emph{shape
analysis}~\cite{ptr_shape}.  It is an analysis of dynamically created
structures like lists. It is able to recognize operations like list creation
and give this information to the analyzer. It can be used for finding memory
leaks or tell whether the list is created properly in the first place.

\paragraph{Field-sensitivity} In languages allowing structured constructs, like
\code{struct} in C, we can either consider each of the member of such an
object when computing the pointer sets or ignore the members. Obviously
ignoring that information leads to a faster analysis, but imprecisions occur.
The imprecisions are two-fold: \begin{enumerate} \item Each member can be
itself a pointer and can have its own pointer target set. Storing this
information only per each structure in field-insensitive sense introduces a
big imprecision.  \item Knowing exactly to which member of a structure a
pointer points is indeed useful information.  \end{enumerate} The
implementation is possible in two ways, depending on the demanded precision.
First, \emph{field-based} sensitivity have only one global information about
structure members, per each member name. So when we have two totally different
structures \code{a} and \code{b} containing a member \code{void *x}, this
analysis considers them to be the same and updates \code{x}'s points-to set
any time \code{a.x} or \code{b.x} is set to point to a new object. The other
analysis, the true \emph{field-sensitive} one, properly handles members of
structures in the natural sense, but for a higher computation price, of
course.

For more information about pointer analysis, see for example Hind's
discussion~\cite{ptr_solved}.

\section{The Goals}
This thesis studies applications of contemporary automatic bug-finding
techniques to large software projects. It specifically aims at techniques like
abstract interpretation or symbolic execution. Both of them were designed in
1970's and were adopted in many tools since then.  The thesis should summarize
the current state of the art of chosen techniques implemented in code
analyzers which can handle large code bases.  On the top of the explored state
of the art, the goal is to develop our techniques which can handle the Linux
kernel.  There should be only few false reports resulting from an
implementation of the techniques.  Another aim is to help others with code
analyzers development. We want to establish a test-bed with predefined
categories of errors usually present in code such as the Linux kernel. For
that purpose a publicly available benchmark should be designed. It should
contain information about errors and false positives found by several
bug-finding tools.  There is also a software competition tailored to code
analyzers. So the goal is to compete with our analyzers.

\section{Thesis Contribution}
We implemented \stan, a framework for code analysis base prevalently on the
ideas of the abstract interpretation. Further we improved the false reports
rate of \stan in \symbio by using a combination of three techniques:
instrumentation, slicing, and symbolic execution. We used \stan also to create
a testbed for new code analysis tools \cdb. Finally, we took \symbio and tuned
it. We used that version on the software competition of code analyzers.

\section{Thesis Organization}
Chapter~\ref{sec:state} deals with a description of techniques and their
current state of the art. In Chapter~\ref{sec:stanse}, we discuss the
abstract interpretation, its implementation in our tool \stan~\cite{stanse}
and its pros and cons. \stan is a framework for developers who wish to
implement some analysis. \stan provides many helpers, e.g.\ for traversing and
checking code. We continue in Chapter~\ref{sec:fmics} by an approach to
eliminate false positives resulting from the abstract interpretation. This is
implemented in the tool called \symbio~\cite{symbiotic}. It consists of
combination of three techniques. First, an instrumentor injects state automata
into the code that is to be checked. Second, a static slicer removes the code
irrelevant to the state machines.  Finally, a symbolic executor evaluates the
code with no false positives on the output.  Chapter~\ref{sec:clabure} then
implements a database of known errors in large projects like the Linux kernel.
The database is called \cdb~\cite{claburedb}. Its intended use is evaluation
of a newly implemented tool. In particular to see how good the reports from
the tool are.  Chapter~\ref{sec:svcomp} shows results of \symbio in a software
competition~\cite{symbiotic-svcomp}.  Finally, we conclude with summaries of
what we have done.
